\pdfoutput=1
\documentclass[11pt]{article}

\usepackage[margin=1in]{geometry}
\usepackage{amsmath}
\usepackage{amssymb}
\usepackage{booktabs}
\usepackage{enumitem}
\usepackage{url}

\title{XScientist: A Long-Running Autonomous Scientific Research System}

\author{
Jixiang Luo\\
Independent Researcher
}

\date{July 2026}

\begin{document}
\maketitle

\begin{abstract}
Autonomous research systems are often evaluated as one-shot paper generators:
given a topic, they produce a manuscript and a small set of experiment logs.
This framing hides the operational problem that makes such systems difficult
to trust: research is long-running, branching, failure-prone, and dependent on
auditable handoffs between agents and humans. XScientist is a research
operating system for this setting. It orchestrates idea generation,
experiment execution, manuscript drafting, self-review, repair, quality
gating, daemon scheduling, and reproducibility artifacts as one continuously
observable pipeline. The central design choice is to treat each run as a
portable research artifact rather than only as a PDF. XScientist exports an
Agent-Native Research Artifact (ARA), a git-like protocol that records an
exploration DAG, per-node code and outputs, claim-to-evidence anchors,
content hashes, provenance, and re-execution hooks. This makes each generated
paper inspectable as a science exploration tree: failed branches, repaired
experiments, ablations, and manuscript claims remain connected to the nodes
that produced them. The system also includes deterministic integrity
forensics, sample gates, truth contracts, reviewer-oriented repair loops, and
long-running daemon controls. This paper describes the current XScientist
architecture, the ARA protocol surface, and the practical safeguards needed
to move autonomous science from single-run demos toward reproducible,
reviewable, and forkable research infrastructure.
\end{abstract}

\section{Introduction}

Large language models have made it possible to automate substantial portions
of the scientific workflow: ideation, experiment scripting, result
interpretation, manuscript drafting, and review simulation. However, a system
that generates a plausible manuscript is not necessarily a system that can be
trusted as research infrastructure. The hard problem is not only text
generation. It is the management of evidence, failures, revisions, and
handoffs over time.

XScientist addresses this operational problem. The system is designed around
four principles:

\begin{enumerate}[leftmargin=*]
  \item \textbf{Research runs should be observable.} Each stage should emit
  structured artifacts that can be inspected without replaying the entire
  workflow.
  \item \textbf{Research runs should be forkable.} A later agent should be
  able to continue from a specific experiment node instead of cold-starting
  from the final PDF.
  \item \textbf{Research claims should be grounded.} Manuscript assertions
  should remain linked to the experiment nodes and evidence artifacts that
  support them.
  \item \textbf{Research automation should expose uncertainty.} Failed
  branches, reviewer objections, integrity warnings, and blocked gates should
  be preserved as first-class outputs rather than hidden as pipeline noise.
\end{enumerate}

The result is a long-running autonomous scientific research system with a
git-like artifact protocol. XScientist can be used as a local project runner,
a continuous paper generator, or a daemon that schedules multiple research
sources over time. Its output is not just a manuscript but a structured
record of how the manuscript came to exist.

\section{System Overview}

XScientist is organized as a research operating system. Figure
\ref{fig:pipeline} summarizes the high-level control flow. The pipeline begins
with topics or source queues, produces ranked ideas, executes experiments,
generates manuscripts, performs self-review and repair, applies quality and
integrity gates, and archives reusable artifacts.

\begin{figure}[t]
\centering
\fbox{
\begin{minipage}{0.92\linewidth}
\centering
\vspace{0.5em}
Topic / Sources
$\rightarrow$ Ideation and Ranking
$\rightarrow$ Experiments
$\rightarrow$ Writeup and Quality Gates
$\rightarrow$ Self-Review and Repair
$\rightarrow$ Artifacts, Index, and Dossier
$\rightarrow$ Daemon Strategy Feedback
\vspace{0.5em}
\end{minipage}
}
\caption{Simplified XScientist control loop. The daemon can feed outcomes
back into later source selection, quality policy, and writing strategy.}
\label{fig:pipeline}
\end{figure}

\subsection{Execution Modes}

The current repository exposes three main entry points.
\begin{itemize}[leftmargin=*]
  \item \texttt{run\_project.py} runs an end-to-end project from a topic or
  idea file.
  \item \texttt{continuous\_paper\_generator.py} performs batch and continuous
  generation across selected ideas and paper types.
  \item \texttt{continuous\_research\_daemon.py} runs a long-lived scheduling
  loop with source queues, dashboarding, reports, and operator controls.
\end{itemize}

\subsection{Pipeline Artifacts}

XScientist writes structured JSON and Markdown artifacts throughout the run:
idea cards, research plans, experiment registries, reviewer findings, repair
plans, manuscript candidates, quality reports, integrity forensics reports,
source provenance, and handoff briefs. These artifacts are deliberately more
important than transient console output. They give downstream agents and
human operators a stable interface for audit and continuation.

\section{Agent-Native Research Artifacts}

The most important interface in XScientist is the Agent-Native Research
Artifact (ARA). An ARA is a directory rooted at a \texttt{manifest.json} file
and an \texttt{exploration\_graph.json} file. It is designed to be read by
another agent without requiring that agent to reverse-engineer the final
paper.

\subsection{Exploration DAG}

Each ARA stores an exploration DAG whose nodes represent concrete experiment,
repair, failure, ablation, or manuscript-candidate states. Edges encode the
parent-to-child evolution of the research process. The system validates that
the graph is directed and acyclic, writes a summary file, and can render a
self-contained HTML visualization. This turns each paper into a science
exploration tree.

\begin{figure}[t]
\centering
\fbox{
\begin{minipage}{0.92\linewidth}
\centering
\vspace{0.5em}
\texttt{root}
$\rightarrow$ \texttt{experiment}
$\rightarrow$ \texttt{failed branch}
$\rightarrow$ \texttt{repair}\\
\texttt{experiment}
$\rightarrow$ \texttt{ablation}
$\rightarrow$ \texttt{candidate paper}
$\rightarrow$ \texttt{claim anchor}
$\rightarrow$ \texttt{fork seed}
\vspace{0.5em}
\end{minipage}
}
\caption{Conceptual science exploration tree. ARA preserves both successful
and failed branches, allowing later agents to audit or fork a specific node.}
\label{fig:tree}
\end{figure}

\subsection{Content Hashing and Provenance}

ARA uses content hashes to identify node payloads and manifest revisions.
Hashes are computed over canonical payloads such as code, metrics, and
optional provenance fields. Forked runs preserve parent pointers, so a child
research run can describe which previous ARA and node seeded it. The protocol
therefore behaves like a lightweight object model for research state.

\subsection{Claim Anchoring}

The writing stage can insert invisible \texttt{\textbackslash claimref}
markers after quantitative or evidence-sensitive claims. A claim registry
scans the manuscript source and writes claim records under the ARA. This
creates a two-way link between paper assertions and the experiment nodes that
support them. Claim coverage can then be summarized and used by quality gates
or downstream reviewers.

\subsection{Fork, Diff, Log, and Verify}

The companion CLI \texttt{run\_ara\_fork.py} exposes operational verbs:
\texttt{inspect}, \texttt{exec}, \texttt{fork}, \texttt{freeze},
\texttt{validate}, \texttt{verify}, \texttt{diff}, \texttt{log}, and
\texttt{refs}. These commands let an operator inspect a node, re-execute its
code, create a new seed ARA, compare two ARAs, or walk the provenance chain.
The same graph data underlies the HTML visualization, the log view, and
node-level diffs.

\section{Quality and Integrity Controls}

Autonomous research systems require checks that are deterministic enough to
be run in CI and explicit enough to be audited later. XScientist implements
several such controls.

\subsection{Self-Review and Repair}

The self-review subsystem runs LLM-based and VLM-assisted reviews, extracts
structured issues, prioritizes repairs, and records repair attempts. The
system tracks whether reviewer objections were addressed, whether manuscript
changes regressed earlier claims, and whether quality thresholds were met.

\subsection{Truth Contracts}

Truth contracts translate a research plan into explicit constraints. They
separate objective facts, comparability constraints, artifact-binding rules,
branch-state rules, and value guardrails. For example, a manuscript should not
claim a baseline-comparable improvement unless the experiment registry
contains a completed run on the declared dataset, metric, and baseline.

\subsection{Sample Gates}

Before full generation, sample gates can require that a small planned task
has a completed experiment record, a budget audit, passing acceptance checks,
and a result summary. This prevents the system from scaling a flawed plan into
a full paper when the cheap preliminary evidence is missing.

\subsection{Deterministic Integrity Forensics}

XScientist adapts an evidence-ledger pattern inspired by
Anti-Autoresearch into the manuscript stage. The integrity forensics pass is
not an authorship detector. It builds a span-anchored ledger from the
manuscript source, applies model-free consistency checks, and writes a
deterministic adjudication report. Hard findings can block submission-ready
status; soft findings are reported for reviewer attention.

\section{Long-Running Daemon Operation}

The daemon mode treats research as an ongoing operation rather than a single
script invocation. It supports source queues, source health tracking,
dashboard serving, daily reports, handoff briefs, pause/resume controls,
source boosting and cooldown, failure protection, and strategy feedback.

This matters because autonomous research systems accumulate operational
state. A topic source may produce weak ideas for a week and then recover. A
review-hardening source may be valuable only after enough experimental
evidence exists. XScientist therefore records source-level and portfolio-level
signals so later cycles can adjust scheduling rather than replay a fixed
script.

\section{Comparison to One-Shot Paper Generation}

Table \ref{tab:comparison} contrasts the XScientist design with a typical
one-shot autonomous paper generator.

\begin{table}[t]
\centering
\begin{tabular}{p{0.29\linewidth}p{0.29\linewidth}p{0.32\linewidth}}
\toprule
Dimension & One-shot generator & XScientist \\
\midrule
Primary output & PDF or draft & PDF plus structured research artifact \\
Failure handling & Often discarded & Preserved as DAG branches \\
Continuation & Prompt-level restart & Node-level fork from ARA \\
Claim grounding & Implicit in prose & Explicit claim-to-node anchors \\
Review & Text feedback & Structured issues, repair attempts, gates \\
Operations & Manual reruns & Daemon, reports, dashboard, controls \\
\bottomrule
\end{tabular}
\caption{Design contrast between one-shot generation and XScientist's
artifact-centered workflow.}
\label{tab:comparison}
\end{table}

\section{Limitations}

XScientist does not make autonomous research automatically correct. The
system can organize evidence, preserve provenance, and expose warnings, but
human review remains necessary for scientific validity, ethical judgment, and
final submission decisions. The current system also has several engineering
limitations:

\begin{itemize}[leftmargin=*]
  \item Re-execution can be expensive or environment-dependent when nodes
  call external APIs, GPUs, or web services.
  \item Integrity forensics is deterministic but incomplete; it catches some
  classes of inconsistency and overclaiming, not all scientific errors.
  \item ARA is a protocol surface rather than a universal standard; broader
  adoption requires independent producers and consumers.
  \item The daemon improves operations, but long-running autonomy still needs
  explicit budgets, stop conditions, and operator oversight.
\end{itemize}

\section{Discussion}

The central claim of XScientist is that autonomous research should be judged
not only by the final manuscript but by the quality of the artifact trail that
produced it. A generated paper without a recoverable experiment tree is hard
to audit. A failed branch without provenance is hard to learn from. A claim
without a node anchor is hard to verify. XScientist makes these relations
first-class.

The same design suggests a path toward collaborative autonomous science:
systems can exchange ARAs, fork one another's nodes, verify claims without
parsing PDFs, and use git-like diffs to review scientific changes. This does
not remove the need for human scientists. It gives human reviewers and
downstream agents a stronger substrate to inspect.

\section{Conclusion}

XScientist is a long-running autonomous scientific research system built
around structured artifacts, continuous operation, and a git-like research
protocol. Its ARA export turns each paper into a science exploration tree with
claim anchors, content hashes, provenance, and re-execution hooks. Together
with self-review, repair loops, sample gates, truth contracts, deterministic
integrity forensics, and daemon scheduling, this design moves autonomous
research from one-shot paper generation toward reproducible and forkable
research infrastructure.

\section*{Availability}

The XScientist repository is available at
\url{https://github.com/smileformylove/XScientist}. The arXiv source for this
system report is maintained under \texttt{paper/xscientist\_arxiv/} in the
repository.

\section*{Acknowledgements}

XScientist builds on ideas from AI Scientist, autoresearch, AIDE,
DeepReviewer, and related open-source autonomous research and review systems.
The integrity forensics layer adapts an evidence-ledger pattern inspired by
Anti-Autoresearch.

\begin{thebibliography}{9}

\bibitem{ai-scientist}
Sakana AI.
\newblock The AI Scientist.
\newblock \url{https://github.com/SakanaAI/AI-Scientist}.

\bibitem{autoresearch}
Andrej Karpathy.
\newblock autoresearch.
\newblock \url{https://github.com/karpathy/autoresearch}.

\bibitem{aide}
WecoAI.
\newblock AIDE.
\newblock \url{https://github.com/WecoAI/aideml}.

\bibitem{deepreviewer}
ResearAI.
\newblock DeepReviewer-v2.
\newblock \url{https://github.com/ResearAI/DeepReviewer-v2}.

\bibitem{anti-autoresearch}
wanshuiyin.
\newblock Anti-Autoresearch.
\newblock \url{https://github.com/wanshuiyin/Anti-Autoresearch}.

\end{thebibliography}

\end{document}
